% titel: Blatt 1, A3: Masse und Bindungsenergie
% Teil II: der durchgerechnete Weg.

\subsection*{Gegeben und gesucht}
${}^{19}\mathrm F$, $Z=9$, $N=10$, Atommasse $m_\text{F}=18.9985\ \mathrm u$.
(a) Kernzusammensetzung; (b) Gesamtmasse der Bestandteile; (c) relativer
Massendefekt; (d) beste Erklärung.

\subsection*{a) Kern}
Protonen $Z=9$, Neutronen $N=A-Z=19-9=10$. Zusammen 19 Nukleonen.

\subsection*{b) Gesamtgewicht der Bestandteile}
Mit $m_\text{p}=1.007276\ \mathrm u$, $m_\text{n}=1.008665\ \mathrm u$,
$m_\text{e}=0.000549\ \mathrm u$ und 9 Elektronen:
\begin{equation}
  m_\text{Teile}=9\,m_\text{p}+10\,m_\text{n}+9\,m_\text{e}
  =19.1571\ \mathrm u .
\end{equation}
(Man zählt das Atom, also Elektronen dazu.)

\subsection*{c) Relativer Massendefekt}
\begin{equation}
  \Delta m = m_\text{Teile}-m_\text{F} = 19.1571 - 18.9985 = 0.1586\ \mathrm u,
\end{equation}
\begin{equation}
  \frac{\Delta m}{m_\text{F}} = \frac{0.1586}{18.9985} \approx 0.00835
  = 0.835\% .
\end{equation}
Die fertige Atommasse ist also etwa $0.84\%$ \emph{leichter} als die Summe
der einzelnen Bestandteile.

\subsection*{d) Erklärung}
Richtig ist: \emph{die fehlende Masse wurde bei der Kernbildung als
Bindungsenergie freigesetzt} ($\Delta E=\Delta m\,c^2$). Die übrigen
Antworten sind Ablenker.

\begin{ergebnis}
$Z{=}9,\ N{=}10$; $m_\text{Teile}=19.157\ \mathrm u$; Massendefekt
$\approx0.1586\ \mathrm u$ bzw.\ $0.84\%$; Begründung über Kernbindungsenergie.
\end{ergebnis}

\begin{kontrolle}
Ein Massendefekt von ca.\ $1\%$ pro Nukleon liegt in der erwarteten
Grössenordnung für leichte Kerne --- plausibel.
\end{kontrolle}

\begin{merksatz}
$E=mc^2$ erklärt, warum Atome leichter sind als ihre Einzelteile: Die beim
Kernaufbau freigesetzte Bindungsenergie erscheint als fehlende Masse.
\end{merksatz}
